\subsection{Signal extraction}
\label{sec:signal-extraction}

We inventory \textbf{unsupervised} candidate features---no $r(q)$ at extraction time---organized into three types.
Each type rests on a \textbf{signal assumption} (why it might predict routing need); hypotheses H1--H3 test these assumptions empirically.

\paragraph{Signal assumptions.}
\begin{itemize}
  \item \textbf{Query-derived:} More complex or ambiguous queries are more likely to require a higher-capability pool member $\Rightarrow$ \textbf{H1}.
  \item \textbf{Model-response:} Higher model uncertainty reflects higher probability of error on that query $\Rightarrow$ \textbf{H2}.
  \item \textbf{Cross-model comparative:} Differences between pool members on the same query reveal routing opportunities $\Rightarrow$ \textbf{H3}.
\end{itemize}
Assumptions are priors, not claims; a rejected hypothesis means the assumption failed for our locked setting.

\paragraph{Query-derived.}
Computed from query text only (no model forward pass): e.g., lexical/syntactic complexity, length, ambiguity cues.

\paragraph{Model-response.}
Computed from one model on $q$: e.g., token entropy $H(q\mid M_i)$, confidence, log-probability, paraphrase stability under rephrasings of the same intent.

\paragraph{Cross-model comparative.}
Computed by comparing pool members on the same $q$, typically from model-response quantities already computed: e.g., entropy gap $\Delta H$, confidence gap, answer disagreement.

Table~\ref{tab:setup} lists concrete signals once the setting is locked (Stages 1--3).
Signal definitions depend on the frozen prompt protocol; implementation details are documented at lock time.

The output of this stage is a raw signal inventory per query---not yet selected into $x(q)$.
